% !Mode:: "TeX:UTF-8"
% 结论页（模板要求的独立结论页，与第6章总结不同）
\chapter*{结论\markboth{结论}{}}
\addcontentsline{toc}{chapter}{结论}

本文构建了一个基于EDA工具反馈的LLM Verilog生成与自动优化原型系统，并通过七组递进式实验验证了反馈循环的有效性和边界条件。

GPT-5.4在RTLLM\_STUDY\_12上的通过率从50\%提升至83\%，4轮重复实验确认了这一结论的稳定性（79.2\%$\pm$4.2\%）。消融实验表明反馈信息本身贡献了约57\%的提升，且存在仅反馈可解而多次独立重试均无法通过的题目。同时，Sonnet 4.6上反馈未能超越纯重试（0/4轮），说明反馈循环的有效性并非对所有模型普适成立。

基于实验结果，有三条对实际部署有参考价值的建议。第一，在目标模型上部署反馈循环前，应先用少量题目做小规模验证，因为Sonnet的结果表明反馈并非对所有模型都有正向效果。第二，反馈粒度建议从L2（仅编译错误）或L3（简洁反馈）开始，过于详尽的反馈反而可能干扰模型的修复判断。第三，对于零样本通过率已较高的简单任务，反馈循环带来的额外收益可能不抵API调用成本的增加；更合理的策略是先尝试零样本生成，仅对失败的题目启用反馈。

本文的工作为理解EDA工具反馈在LLM RTL自动生成中的作用提供了实证分析和工程参考。
